\section{The exact two-loop anomaly}\label{sec:two-loop}

\subsection{The two geometric integrals}

A connected cubic graph with geometric loop number $n$ contributes a
local contact only if
\[
 2|E|=3|V|-4,\qquad |E|=3n+1,\qquad |V|=2n+2 .
\]
It then has four external fields. At two loops the two surviving
graphs are theta graphs whose three paths have lengths $(1,3,3)$
and $(2,2,3)$. The four external fields occupy their bivalent
vertices. Each graph has automorphism order four.

For a theta graph with path lengths $n_1,n_2,n_3$, the integration
over its path Schwinger parameters gives, in the present dimension,
the following expression in the Euler gamma function $\Gamma$:
\[
 \frac12
 \frac{\Gamma(\frac72-n_1)\Gamma(\frac72-n_2)
       \Gamma(\frac72-n_3)}
      {\Gamma(\frac72)(n_1-1)!(n_2-1)!(n_3-1)!}.
\]
The oriented momentum determinant supplies the graph-dependent sign.
In the orientation
of the original graph tensors the complete results are
\begin{align}
 K_{133}
 &=-\frac12\frac{\Gamma(\frac52)\Gamma(\frac12)^2}
                    {4\Gamma(\frac72)}
       [p_0p_1][p_2p_3]
   =-\frac{\pi}{20}[p_0p_1][p_2p_3],\notag\\
 K_{223}
 &=\frac12\frac{\Gamma(\frac32)^2\Gamma(\frac12)}
                    {2\Gamma(\frac72)}
       [p_0p_1][p_2p_3]
   =\frac{\pi}{30}[p_0p_1][p_2p_3].
 \label{eq:theta-kernels}
\end{align}
These are projective Schwinger integrals at zero external
displacement, with arbitrary conserved holomorphic momenta.
Their derivation by successive path integration and their original
orientation certificates are given in \cite{RepoOriginal}.

Let $W_{133}$ and $W_{223}$ denote the color contractions of the
six classical cubic tensors at the vertices, with $\kappa^{-1}$
on the seven edges, including the external-field ordering sign.
The subscript records the three path lengths. Define their full
alternating Taylor cochains by
\begin{equation}
 F_{133}=\alt\bigl([p_0p_1][p_2p_3]W_{133}\bigr),\qquad
 F_{223}=\alt\bigl([p_0p_1][p_2p_3]W_{223}\bigr).
 \label{eq:theta-cochains}
\end{equation}
With $\nu=-U/(64\pi^5)$, the automorphism factors and
\eqref{eq:theta-kernels} give
\begin{equation}
 \boxed{\quad
 o_2=\frac{\pi\nu}{240}(-3F_{133}+2F_{223})
     =-\frac{U}{15360\pi^4}(-3F_{133}+2F_{223}).
 \quad}
 \label{eq:o2}
\end{equation}
In particular the relative coefficients are $-3$ and $2$;
the scalar $U$ is common to the two graphs.

\subsection{Full color tensors for traceless matrices}

For $\g=\fsl_N$ with $\kappa=\tr$, the original contractions have
the following complete expressions:
\begin{align}
 W_{133}
 &=2\tr(X_0X_1X_2X_3)
   +(2+N^2)\tr(X_0X_1X_3X_2)\notag\\
 &\quad-4\tr(X_0X_2X_1X_3)
   +(2+N^2)\tr(X_0X_2X_3X_1)\notag\\
 &\quad-4\tr(X_0X_3X_1X_2)
   +2\tr(X_0X_3X_2X_1)
   +6N\tr(X_0X_1)\tr(X_2X_3),\notag\\
 W_{223}
 &=-2\tr(X_0X_1X_2X_3)-2\tr(X_0X_1X_3X_2)\notag\\
 &\quad+4\tr(X_0X_2X_1X_3)-2\tr(X_0X_2X_3X_1)\notag\\
 &\quad+4\tr(X_0X_3X_1X_2)-2\tr(X_0X_3X_2X_1)
       -4N\tr(X_0X_1)\tr(X_2X_3)\notag\\
 &\quad+2N\tr(X_0X_2)\tr(X_1X_3)
       +2N\tr(X_0X_3)\tr(X_1X_2).
 \label{eq:complete-theta-colors}
\end{align}
No choice of external matrices or momenta is imposed here.
Write
\[
 f(X_0,X_1,X_2)=\kappa(X_0,[X_1,X_2]),\qquad
 \beta=d\bigl([p_0p_1]^2f\bigr).
\]
In the trace normalization, full alternation of
\eqref{eq:complete-theta-colors} gives the cochain identity
\begin{equation}
 F_{133}+F_{223}=2N^2\beta,\qquad
 o_2=\pi\nu\left(\frac{F_{223}}{48}
                         -\frac{N^2\beta}{40}\right).
 \label{eq:sl-theta-boundary}
\end{equation}
Thus the bare obstruction class is
\[
 [o_2]=-\frac{U}{3072\pi^4}[F_{223}].
\]
Classical closedness follows either directly from these graph
cochains or from the first unsolved coefficient of the
regularized master equation, since the one-loop contact vanishes.

\subsection{A detector valid for every simple Lie algebra}

Use the Killing form in this subsection. For an invariant
four-cochain of momentum degree four, take the coefficient of
$[p_0p_1][p_0p_2]$. Denote its full color contraction with
$\kappa(X_0,X_1)\kappa(X_2,X_3)$ by $\ell_{kk}$, and its
contraction with $\kappa([X_0,X_1],[X_2,X_3])$ by $\ell_{ff}$.
Every contraction uses the inverse Killing form. The three
determinants $[p_0p_1],[p_0p_2],[p_1p_2]$ are algebraically
independent, so this coefficient is unambiguous.

The boundary $\beta$ has
\[
 \ell_{kk}(\beta)=2\dim\g,\qquad
 \ell_{ff}(\beta)=\dim\g.
\]
Consequently the linear functional
\[
 \varpi=(\dim\g)\ell_{kk}-2(\dim\g)\ell_{ff}
 \qquad\text{satisfies}\qquad \varpi(\beta)=0 .
\]
It annihilates every boundary in the invariant sector in question.
Indeed exactness of invariants for the finite-dimensional
representations of the adjoint group and $SL_2$ at fixed momentum degree supplies
an invariant primitive whenever an invariant cochain is exact.
A cubic primitive of degree four has momentum polynomial
$[p_0p_1]^2$ and an invariant alternating color tensor.
That tensor is a multiple of $f$ for every simple algebra.
Thus the entire invariant boundary space is $\mathbb C\beta$.

To evaluate the detector, let $e_b,e^b$ be Killing-dual bases,
let $P$ exchange the two factors of $\g\otimes\g$, and put
\[
 \Omega_{\mathrm{ad}}
   =\sum_b\operatorname{ad}(e_b)\otimes\operatorname{ad}(e^b),
 \qquad n_4=\Tr_{\g\otimes\g}(P\Omega_{\mathrm{ad}}^4).
\]
Contraction of the original theta networks gives
\begin{align}
 \ell_{kk}(F_{223})&=-6\dim\g,&
 \ell_{ff}(F_{223})&=-8n_4,\notag\\
 \varpi(F_{223})&=2\dim\g(8n_4-3\dim\g),&
 \varpi(F_{133})&=-\varpi(F_{223}).
 \label{eq:universal-detector}
\end{align}
For clarity, these contractions can be reduced to the three
pairings $(01)(23),(02)(13),(03)(12)$ before alternation.
For $W_{223}$ their metric contractions are
$(-\dim\g/2,\dim\g/4,\dim\g/4)$ and their two-bracket contractions
are $(0,n_4,n_4)$. Full alternation applies the linear combinations
$(8,-4,-4)$ and $(0,-4,-4)$, respectively, giving
\eqref{eq:universal-detector}.

Here is the exact family-independent evaluation of the remaining
Casimir trace. In Vogel's normalization let
$(v_1,v_2,v_3)$ be the three parameters and set
$t_V=v_1+v_2+v_3$. The fourth adjoint Casimir formula
\cite{MSV} gives
\begin{align}
 n_4&=\dim\g\left[
  \frac{3t_V(v_1^2+v_2^2+v_3^2)
              -(v_1^3+v_2^3+v_3^3)}
       {16t_V^3}-\frac18\right],\notag\\
 8n_4-3\dim\g
 &=-\frac{3\dim\g\,(t_V+v_1)(t_V+v_2)(t_V+v_3)}
             {2t_V^3}.
 \label{eq:vogel-detector}
\end{align}
The substitutions and results are displayed together:
\begin{center}
\renewcommand{\arraystretch}{1.55}
\begin{tabular}{@{}lll@{}}
\toprule
$\g$ & $(v_1,v_2,v_3)$ & $8n_4-3\dim\g$\\
\midrule
$\fsl_N$ & $(-2,2,N)$
  & $-\dfrac{3(N^2-1)(N^2-4)}{N^2}$\\
\addlinespace[3pt]
$\fso_N$ & $(-2,4,N-4)$
  & $-\dfrac{3N(N-1)(N-3)(N-4)(N+2)}{2(N-2)^3}$\\
\addlinespace[3pt]
$\fsp_N$ & $(-2,1,\frac{N+4}{2})$
  & $-\dfrac{3N(N-2)(N+1)(N+3)(N+4)}{2(N+2)^3}$\\
\addlinespace[3pt]
$G_2$ & $(-2,\frac{10}{3},\frac83)$ & $-\dfrac{385}{12}$\\
\addlinespace[3pt]
$F_4$ & $(-2,5,6)$ & $-\dfrac{12740}{81}$\\
\addlinespace[3pt]
$E_6$ & $(-2,6,8)$ & $-\dfrac{975}{4}$\\
\addlinespace[3pt]
$E_7$ & $(-2,8,12)$ & $-\dfrac{34580}{81}$\\
\addlinespace[3pt]
$E_8$ & $(-2,12,20)$ & $-\dfrac{12152}{15}$\\
\bottomrule
\end{tabular}
\end{center}
For $\fsp_N$ the defining representation has even dimension $N$.
The zeros in the simple families occur precisely at
$\fsl_2=\fso_3=\fsp_2$; $\fso_4$ is not simple.

\begin{theorem}[Nontriviality of the bare anomaly]
\label{thm:bare-anomaly}
The bare two-loop anomaly is a nonzero cohomology class for every
complex simple gauge algebra other than $\fsl_2$.
\end{theorem}
\begin{proof}
Equations \eqref{eq:o2} and \eqref{eq:universal-detector} give
\[
 \varpi(o_2)=\frac{\pi\nu}{24}
                 \dim\g(8n_4-3\dim\g).
\]
The table makes this nonzero in the stated cases, whereas
$\varpi$ vanishes on every possible invariant boundary.
In Killing normalization for $\fsl_2$, direct contraction gives
$F_{133}=\frac72\beta$ and $F_{223}=-3\beta$; hence its
obstruction is exact. Rescaling the invariant form by $c\ne0$
multiplies each six-vertex, seven-edge theta graph by
$c^{6-7}=c^{-1}$ and cannot change nontriviality.
\end{proof}
